\documentclass[11pt]{article}
\usepackage[margin=1in]{geometry}
\usepackage{amsmath,amssymb,amsthm}
\usepackage{graphicx}
\usepackage[hidelinks]{hyperref}

\newtheorem{theorem}{Theorem}
\newtheorem{proposition}[theorem]{Proposition}
\newtheorem{corollary}[theorem]{Corollary}
\newcommand{\Lip}{\operatorname{Lip}}
\newcommand{\Val}{\operatorname{Val}}
\newcommand{\id}{\operatorname{id}}

\title{Norm continuity of the $\mathrm{GL}(n)$ action\\in low degree and in dimension two}
\author{A substantial partial result for arXiv:2401.05913}
\date{11 August 2026}

\begin{document}
\maketitle

\begin{center}
\fbox{\parbox{0.91\textwidth}{\textbf{Status: substantial partial result,
likely valid.}  The action is norm-continuous on all valuations when $n=2$
and on the degree-zero and degree-one summands in every dimension.  For
$n\geq3$, only a possibly nonsmooth degree-two summand remains.}}
\end{center}

\section{Source question and notation}

Let $S^{n-1}$ be the Euclidean unit sphere and equip $\Lip(S^{n-1})$ with
\[
 \|f\|_{\Lip}=\max\{\|f\|_\infty,\operatorname{lip}(f)\}.
\]
Write $\mathcal V_n$ for the space of complex-valued, $\tau$-continuous,
dually translation invariant valuations on $\Lip(S^{n-1})$.  Thus
$\mu(f+\ell)=\mu(f)$ whenever $\ell$ is the restriction of a linear
functional on $\mathbb R^n$.  The source \cite{CKP} equips this space with
the Banach norm
\[
 \|\mu\|=\sup_{\|f\|_{\Lip}\leq1}|\mu(f)|
\]
and, in Remark~4.3 on PDF page~18, asks whether the natural action
\[
 [\pi(g)\mu](f)=\mu(g^{-1}\mathbin{\cdot}f),\qquad g\in\mathrm{GL}(n,\mathbb R),
\]
is continuous for this norm.  Here
\[
 (g\mathbin{\cdot}f)(x)=|g^Tx|f\left(\frac{g^Tx}{|g^Tx|}\right).
\]

\begin{figure}[ht]
\centering
\includegraphics[width=0.95\textwidth]{figures/source_question.png}
\caption{The norm-continuity question in \cite[Remark~4.3]{CKP}.}
\end{figure}

The source proves the homogeneous decomposition
\[
 \mathcal V_n=\mathcal V_{n,0}\oplus\mathcal V_{n,1}
 \oplus\mathcal V_{n,2}.
\]
Our result settles the first two summands and all three when $n=2$.

\section{Idea of the proof}

Restriction to support functions gives an injective equivariant map from the
functional valuations to classical translation invariant valuations on convex
bodies.  In degree one, McMullen polynomiality says that the classical
valuation is Minkowski additive, hence integration of the support function
against a finite measure.  The same integral defines a valuation on every
Lipschitz function, and injectivity forces it to be the original valuation.
Small linear changes move a Lipschitz function uniformly by a small amount,
uniformly over the unit ball, so integration against the fixed measure gives
norm-continuity.

In dimension two the degree-zero part is constant and the classical
degree-two space is spanned by area.  Its known functional extension
$\int(f^2-|\nabla f|^2)$ is therefore the only possible degree-two part, and
it transforms by the continuous scalar $|\det g|^{-1}$.  Local operator
bounds then turn continuity of fixed orbits into joint continuity of the
action.

\section{The low-degree theorem}

Let $T:\mathcal V_n\to\Val(\mathbb R^n)$ be the source's injective map
\[
 T\mu(K)=\mu(h_K),
\]
where $h_K$ is the support function of the convex body $K$.

\begin{theorem}[low degree and dimension two]\label{thm:main}
Let $n\geq2$.
\begin{enumerate}
\item The $\mathrm{GL}(n,\mathbb R)$ action is jointly continuous in the norm
of Remark~4.3 on $\mathcal V_{n,0}\oplus\mathcal V_{n,1}$.
\item Every $\mu\in\mathcal V_{n,1}$ has the form
\[
 \mu(f)=\int_{S^{n-1}}f\,d\sigma,
\]
where $\sigma$ is a finite complex Borel measure satisfying
$\int_{S^{n-1}}x\,d\sigma(x)=0$.
\item For $n=2$, the action is jointly norm-continuous on all of
$\mathcal V_2$.
\end{enumerate}
\end{theorem}

\begin{proof}
A zero-homogeneous continuous valuation is constant: for $t>0$,
$\mu_0(tf)=\mu_0(f)$, and letting $t\downarrow0$ gives
$\mu_0(f)=\mu_0(0)$.  Its orbit is therefore fixed.

Now let $\mu\in\mathcal V_{n,1}$.  The classical valuation $T\mu$ is
continuous, translation invariant, and one-homogeneous.  The degree-one case
of McMullen polynomiality makes it Minkowski additive \cite{McMullen,Ludwig}.
The standard support-function representation therefore gives a finite complex
Borel measure $\sigma$ on $S^{n-1}$ such that
\[
 T\mu(K)=\int_{S^{n-1}}h_K\,d\sigma.
\]
Translation invariance is exactly the zero-first-moment condition on
$\sigma$.  Define $\mu_\sigma(f)=\int f\,d\sigma$.  This is a linear, hence
one-homogeneous, valuation; it is dually translation invariant by the moment
condition and $\tau$-continuous because $\tau$-convergence includes uniform
convergence.  Moreover, $T\mu_\sigma=T\mu$.  Injectivity of $T$ from
\cite[Proposition~4.4]{CKP} yields $\mu=\mu_\sigma$.  (For complex-valued
valuations, apply the real representation to the real and imaginary parts.)

We prove continuity of this orbit at the identity.  Put $h=g^{-1}$ and
\[
 \Phi_h(x)=\frac{h^Tx}{|h^Tx|}.
\]
Uniformly over $x\in S^{n-1}$ and $\|f\|_{\Lip}\leq1$,
\begin{align*}
 |(h\mathbin{\cdot}f)(x)-f(x)|
 &\leq \big||h^Tx|-1\big|\,|f(\Phi_h(x))|
       +|f(\Phi_h(x))-f(x)|\\
 &\leq \|h-I\|_{\rm op}+|\Phi_h(x)-x|.
\end{align*}
The right side tends uniformly to zero as $g\to I$.  Consequently,
\[
 \|\pi(g)\mu-\mu\|
 \leq \|\sigma\|_{\rm TV}
 \sup_{\|f\|_{\Lip}\leq1}\|h\mathbin{\cdot}f-f\|_\infty
 \longrightarrow0.
\]

Suppose next that $n=2$.  The map $T$ embeds
$\mathcal V_{2,2}$ into the one-dimensional classical space
$\Val_2(\mathbb R^2)$ spanned by area.  The nonzero valuation
\[
 Q(f)=\int_{S^1}\big(f(x)^2-|\nabla f(x)|^2\big)\,d\mathcal H^1(x)
\]
belongs to $\mathcal V_{2,2}$ and restricts to a nonzero multiple of area
\cite{CPTV,CKP}.  Hence $\mathcal V_{2,2}=\mathbb C Q$.  Equivariance and
injectivity of $T$ give
\[
 \pi(g)Q=|\det g|^{-1}Q,
\]
so this final orbit is norm-continuous.

For completeness, the representation operators are locally bounded.  The
source's gradient covariance implies
\[
 \|g\mathbin{\cdot}f\|_{\Lip}
 \leq2\|g\|_{\rm op}\|f\|_{\Lip}.
\]
The homogeneous projections are bounded as well: writing
$A=\mu(f/2)-\mu(0)$ and $B=\mu(f)-\mu(0)$ gives
\[
\mu_2(f)=2B-4A,\qquad \mu_1(f)=4A-B,
\]
and therefore $\|\mu_k\|\leq12\|\mu\|$ (a deliberately uniform bound for
$k=0,1,2$).  Positive homogeneity now gives
locally uniform bounds for $\pi(g)$.  Continuity of every fixed orbit proved
above consequently implies joint continuity.
\end{proof}

\section{An integral degree-two extension}

The same mechanism covers a broad part of the remaining degree-two space.
Write
\[
 \bar\nabla f(x)=\nabla f(x)+f(x)x.
\]

\begin{proposition}[integral degree-two class]\label{prop:l1}
Let $\varphi\in L^1(S^{n-1};\operatorname{Sym}^2(\mathbb C^n))$ satisfy
\[
 \int_{S^{n-1}}\varphi(x)\bar\nabla f(x)\,d\mathcal H^{n-1}(x)=0
 \qquad(f\in\Lip(S^{n-1})).
\]
Then
\[
 \mu_\varphi(f)=\int_{S^{n-1}}
 \langle\varphi(x)\bar\nabla f(x),\bar\nabla f(x)\rangle
 \,d\mathcal H^{n-1}(x)
\]
belongs to $\mathcal V_{n,2}$, and its $\mathrm{GL}(n,\mathbb R)$ orbit is
norm-continuous.  The same holds for the norm closure of all such valuations.
\end{proposition}

\begin{proof}
At almost every point, the unordered pair of first jets of $f\vee h$ and
$f\wedge h$ equals that of $f$ and $h$; on the equality set their gradients
agree almost everywhere.  This proves the valuation identity.  Homogeneity
is immediate.  If $\ell_a(x)=\langle a,x\rangle$, then
$\bar\nabla\ell_a=a$.  Expanding $\mu_\varphi(f+\ell_a)$ and using the stated
condition, first for $f$ and then for linear functions, proves dual
translation invariance.  For a $\tau$-convergent sequence the augmented
gradients converge almost everywhere and are uniformly bounded; dominated
convergence against the fixed $L^1$ coefficient proves $\tau$-continuity.  Also
\[
 \|\mu_\varphi\|\leq4\|\varphi\|_{L^1}.
\]

For $h=g^{-1}$, the source's covariance formula reads
\[
 \bar\nabla(h\mathbin{\cdot}f)(x)
 =h\,\bar\nabla f(\Phi_h(x)).
\]
Changing variables $y=\Phi_h(x)$ rewrites
$\pi(g)\mu_\varphi$ as $\mu_{\varphi_g}$, where, with $J_g$ the spherical
Jacobian of $\Phi_g$,
\[
 \varphi_g(y)=J_g(y)h^T\varphi(\Phi_g(y))h.
\]
As $g\to I$, the matrix factors and $J_g$ converge uniformly and pullback by
$\Phi_g$ is strongly continuous on $L^1$.  Thus
$\|\varphi_g-\varphi\|_1\to0$ and
\[
 \|\pi(g)\mu_\varphi-\mu_\varphi\|
 \leq4\|\varphi_g-\varphi\|_1\longrightarrow0.
\]
Local boundedness of the representation passes this conclusion to the norm
closure.
\end{proof}

Every smooth degree-two valuation classified in \cite[Theorem~3]{CKP} has
precisely the form in Proposition~\ref{prop:l1} with a smooth matrix field.
The proposition
therefore includes the norm closure of the source's smooth degree-two
subspace.

\section{Limitations and novelty}

For $n\geq3$, the full question is now confined to degree two outside the
norm closure of the integral class.  The source proves density of smooth
valuations only for uniform convergence on compact subsets of
$\Lip(S^{n-1})$; its unit ball is not compact in the $\tau$ metric because it
contains arbitrarily oscillatory gradients.  Convolution by a group
approximate identity cannot repair this directly: norm convergence of that
convolution is equivalent to the orbit continuity one is trying to prove.

The most obvious negative construction also fails.  A finitely additive or
singular coefficient can make translations discontinuous in a dual norm, but
$\tau$-continuity excludes finitely additive gradient functionals
\cite{CPTV}; singular quadratic coefficients would require evaluating an
almost-everywhere gradient on a null set.  Thus neither a full representation
nor a valid discontinuous orbit has been obtained for the residual summand.

A bounded novelty search through 11 August 2026 checked the run's exact-id and
keyword indexes, exact searches for the Remark~4.3 wording and for
``GL action norm continuity'' with the paper's title and terminology, and the
OpenAlex record for arXiv:2401.05913 (zero indexed citations).  No later answer
or matching low-degree theorem was located.  Novelty confidence is moderate:
the measure argument is classical once the source's injectivity is noticed,
but no prior statement of this consequence was found.

\paragraph{Human-review recommendation.}
Check especially the use of the classical degree-one measure representation,
the passage from injectivity of $T$ to the formula for all Lipschitz
functions, and the $L^1$ change-of-variables formula in
Proposition~\ref{prop:l1}.

\begin{thebibliography}{9}
\bibitem{CKP}
A.~Colesanti, J.~Knoerr, and D.~Pagnini,
\emph{The homogeneous decomposition of dually translation invariant
valuations on Lipschitz functions on the sphere},
arXiv:\href{https://arxiv.org/abs/2401.05913}{2401.05913} (2024).

\bibitem{McMullen}
P.~McMullen,
\emph{Continuous translation invariant valuations on the space of compact
convex sets},
Arch. Math. 34 (1980), 377--384.

\bibitem{Ludwig}
M.~Ludwig and F.~Mussnig,
\emph{Valuations on convex bodies and functions},
arXiv:\href{https://arxiv.org/abs/2302.00416}{2302.00416} (2023),
Corollary~4.9.

\bibitem{CPTV}
A.~Colesanti, D.~Pagnini, P.~Tradacete, and I.~Villanueva,
\emph{Continuous valuations on the space of Lipschitz functions on the
sphere},
J. Funct. Anal. 280 (2021), 108873;
arXiv:\href{https://arxiv.org/abs/2005.05419}{2005.05419}.
\end{thebibliography}

\end{document}
